\begin{theorem}[Wild quadratic refinement identities]\label{mod:cases-wildquadratic-quadraticrefinement}
Let $k$ be the common finite residue field in the wild quadratic
calculation and put
\[
 \psi_0(x)=(-1)^{\operatorname{Tr}_{k/\mathbf F_2}(x)}.
\]
For a normalized characteristic-two refinement $\mathfrak q$, the
polar sign and Frobenius normalization used here are exactly
\[
 \mathfrak q(x+y)=\mathfrak q(x)\mathfrak q(y)\psi_0(xy),
 \qquad \mathfrak q(x^2)=\mathfrak q(x).
\]
Thus the polar character is positive.  Every same-polar critical
function is uniquely
\[
 h_\lambda(z)=\mathfrak q(z)\psi_0(\lambda z).
\]
The four names used later are
$h_{\gamma_0}$ for $\tau$, $h_\gamma$ for $\chi_F$,
$h_{\gamma_1}$ for $\tau\chi_F$, and $h_{\gamma'}$ for $\chi_K$;
the name is used only in a parity row in which that odd critical
function exists.

Every positive-polar Hasse function admits a normalized affine twist.
Relative to a fixed normalized $\mathfrak q$, the two and only two
normalized choices are
\[
 \mathfrak q,\qquad
 \mathfrak q_1(z)=\mathfrak q(z)\psi_0(z),
\]
and they are distinct.  Changing the normalized base changes the
coefficient in the precise direction
\[
 h^{\mathfrak q_1}_{\lambda+1}=h^{\mathfrak q}_\lambda.
\]
The corresponding correction values and phases agree after this
paired shift; the two normalized refinements themselves are never
identified.

The complete residual correction value is
\[
 C_{\mathfrak q}(\lambda,c)
 =h_\lambda(c)
 =\mathfrak q(c)\psi_0(\lambda c)
 =\mathfrak q(\lambda)^{-1}\mathfrak q(c+\lambda).
\]
It is this value, not its inverse or conjugate.  With
\[
 H_{\mathfrak q}(\lambda)
 =\Aphase\!\left(\sum_{z\in k}h_\lambda(z)\right),
\]
the inverse occurs only on the translated phase side:
\[
 H_{\mathfrak q}(\lambda+c)
 =C_{\mathfrak q}(\lambda,c)^{-1}H_{\mathfrak q}(\lambda),
 \qquad
 C_{\mathfrak q}(\lambda,c)H_{\mathfrak q}(\lambda+c)
 =H_{\mathfrak q}(\lambda).
\]

Precomposition by Frobenius sends the affine coefficient to inverse
Frobenius.  In particular, if $(\gamma')^2=L$, then
\[
 h_L(z^2)=h_{\gamma'}(z),
 \qquad H_{\mathfrak q}(\gamma')=H_{\mathfrak q}(L).
\]
The square root is unique in characteristic two.  The remaining
normalization identities exposed downstream are
\[
 \begin{gathered}
 H_{\mathfrak q}(\lambda)
   =H_{\mathfrak q}(0)\overline{\mathfrak q(\lambda)},
 \qquad H_{\mathfrak q}(0)^2=\mathfrak q(1),\\
 H_{\mathfrak q}(\lambda)^2
   =\mathfrak q(1)\psi_0(\lambda),
 \qquad H_{\mathfrak q}(\lambda^2)=H_{\mathfrak q}(\lambda),\\
 H_{\mathfrak q}(1+\lambda)H_{\mathfrak q}(\lambda)=1.
 \end{gathered}
\]

For the boundary coordinate, let $r^2=\rho$ and $1+\rho\ne0$.
Lean records the two exact affine identities
\[
 h_{\gamma_0}(z)h_\gamma(\rho z)
 =h_{(\gamma_0+\rho\gamma+r)/(1+\rho)}((1+\rho)z)
\]
and
\[
 h_\gamma(z)h_{\gamma_0}(\rho z)\psi_0(\rho z)
 =h_{(\gamma+\rho\gamma_0+\rho+r)/(1+\rho)}((1+\rho)z).
\]
Thus all affine coefficient directions, including the terms $+r$ and
$+\rho$, are explicit.  The raw upper functions away from the boundary
are respectively $h_\gamma$ and $h_{1+\gamma_0}$.

Finally, if $c+d=a+b+A$, then the signed phase quotient is
\[
 \frac{H(a)H(b)}{H(c)H(d)}
 =\mathfrak q(A)\psi_0\bigl((a+b)A+cd+ab\bigr).
\]
If also $B+C=A$, multiplication by the positively oriented values
$\mathfrak q(B)\mathfrak q(C)$ leaves
\[
 \psi_0\bigl((a+b)A+cd+ab+A+BC\bigr).
\]
When this exponent plus $t$ equals $z+z^2$, the exact
Artin--Schreier identity $\psi_0(z+z^2)=1$ gives
\[
 H(a)H(b)\mathfrak q(B)\mathfrak q(C)\psi_0(t)=H(c)H(d).
\]
\LeanFile{LanglandsFirstMainLemma/Cases/WildQuadratic/QuadraticRefinement.lean}
\lean{LanglandsFirstMainLemma.WildQuadraticRefinement.criticalFunction}
\lean{LanglandsFirstMainLemma.WildQuadraticRefinement.criticalFunction_apply}
\lean{LanglandsFirstMainLemma.WildQuadraticRefinement.tauCriticalFunction}
\lean{LanglandsFirstMainLemma.WildQuadraticRefinement.chiCriticalFunction}
\lean{LanglandsFirstMainLemma.WildQuadraticRefinement.tauChiCriticalFunction}
\lean{LanglandsFirstMainLemma.WildQuadraticRefinement.chiKCriticalFunction}
\lean{LanglandsFirstMainLemma.WildQuadraticRefinement.tauCriticalFunction_apply}
\lean{LanglandsFirstMainLemma.WildQuadraticRefinement.chiCriticalFunction_apply}
\lean{LanglandsFirstMainLemma.WildQuadraticRefinement.tauChiCriticalFunction_apply}
\lean{LanglandsFirstMainLemma.WildQuadraticRefinement.chiKCriticalFunction_apply}
\lean{LanglandsFirstMainLemma.WildQuadraticRefinement.exists_normalized}
\lean{LanglandsFirstMainLemma.WildQuadraticRefinement.positive_polar}
\lean{LanglandsFirstMainLemma.WildQuadraticRefinement.frobenius_normalized}
\lean{LanglandsFirstMainLemma.WildQuadraticRefinement.existsUnique_criticalFunction}
\lean{LanglandsFirstMainLemma.WildQuadraticRefinement.criticalFunction_frobenius}
\lean{LanglandsFirstMainLemma.WildQuadraticRefinement.criticalFunction_frobenius_of_sq}
\lean{LanglandsFirstMainLemma.WildQuadraticRefinement.squareRoot}
\lean{LanglandsFirstMainLemma.WildQuadraticRefinement.squareRoot_sq}
\lean{LanglandsFirstMainLemma.WildQuadraticRefinement.squareRoot_unique}
\lean{LanglandsFirstMainLemma.WildQuadraticRefinement.secondNormalized}
\lean{LanglandsFirstMainLemma.WildQuadraticRefinement.secondNormalized_apply}
\lean{LanglandsFirstMainLemma.WildQuadraticRefinement.normalized_eq_or_eq_second}
\lean{LanglandsFirstMainLemma.WildQuadraticRefinement.secondNormalized_ne}
\lean{LanglandsFirstMainLemma.WildQuadraticRefinement.secondNormalized_criticalFunction}
\lean{LanglandsFirstMainLemma.WildQuadraticRefinement.correctionValue}
\lean{LanglandsFirstMainLemma.WildQuadraticRefinement.correctionValue_apply}
\lean{LanglandsFirstMainLemma.WildQuadraticRefinement.correctionValue_eq_criticalFunction}
\lean{LanglandsFirstMainLemma.WildQuadraticRefinement.correctionValue_eq_inv_mul}
\lean{LanglandsFirstMainLemma.WildQuadraticRefinement.secondNormalized_correctionValue}
\lean{LanglandsFirstMainLemma.WildQuadraticRefinement.phase}
\lean{LanglandsFirstMainLemma.WildQuadraticRefinement.phase_add}
\lean{LanglandsFirstMainLemma.WildQuadraticRefinement.correction_phase_translation}
\lean{LanglandsFirstMainLemma.WildQuadraticRefinement.secondNormalized_phase}
\lean{LanglandsFirstMainLemma.WildQuadraticRefinement.phase_eq}
\lean{LanglandsFirstMainLemma.WildQuadraticRefinement.phase_zero_sq}
\lean{LanglandsFirstMainLemma.WildQuadraticRefinement.phase_sq}
\lean{LanglandsFirstMainLemma.WildQuadraticRefinement.phase_frobenius}
\lean{LanglandsFirstMainLemma.WildQuadraticRefinement.phase_eq_of_sq}
\lean{LanglandsFirstMainLemma.WildQuadraticRefinement.phase_pair}
\lean{LanglandsFirstMainLemma.WildQuadraticRefinement.absoluteTraceChar_mul_sq_of_sq}
\lean{LanglandsFirstMainLemma.WildQuadraticRefinement.value_mul_scale_root}
\lean{LanglandsFirstMainLemma.WildQuadraticRefinement.rawUpperBelow}
\lean{LanglandsFirstMainLemma.WildQuadraticRefinement.rawUpperAbove}
\lean{LanglandsFirstMainLemma.WildQuadraticRefinement.boundaryLowerProduct}
\lean{LanglandsFirstMainLemma.WildQuadraticRefinement.rawUpperBoundary}
\lean{LanglandsFirstMainLemma.WildQuadraticRefinement.boundaryLowerProduct_eq_affine}
\lean{LanglandsFirstMainLemma.WildQuadraticRefinement.rawUpperBoundary_eq_affine}
\lean{LanglandsFirstMainLemma.WildQuadraticRefinement.phase_four}
\lean{LanglandsFirstMainLemma.WildQuadraticRefinement.phase_four_correction}
\lean{LanglandsFirstMainLemma.WildQuadraticRefinement.artinSchreier_sign}
\lean{LanglandsFirstMainLemma.WildQuadraticRefinement.phase_four_cancel}
\lean{LanglandsFirstMainLemma.wildQuadratic_refinement}
\leanok
\SourceText{Lemmas lem:normalized-refinement and lem:quadratic-correction-rule.}
\uses{mod:finitefield-chartworefinement,mod:finitefield-artinschreier}
\FormalizationContract
\end{theorem}
\begin{proof}
Existence and uniqueness of the normalized affine family, the positive
polar identity, the two normalized choices, correction translation,
phase identities, and signed four-phase formulas are specialized
directly from the certified characteristic-two refinement theory.
The paired coefficient shift is obtained by expanding
$\mathfrak q_1=\mathfrak q\psi_0$ and using
$\psi_0(x)^2=1$.

For the boundary formula, the refinement law first gives
\[
 \mathfrak q(z)\mathfrak q(\rho z)
 =\mathfrak q((1+\rho)z)\psi_0(\rho z^2).
\]
If $r^2=\rho$, Frobenius invariance of the absolute trace changes the
last phase to $\psi_0(rz)$.  Combining it with the two affine phases
and clearing the proved nonzero denominator $1+\rho$ gives the two
displayed coefficients.  Frobenius precomposition uses the unique
square root of a raw coefficient.  The final cancellation is the
certified four-phase formula followed by the literal plus-sign
Artin--Schreier identity.
\end{proof}
